\section{Introduction}
\label{sec:intro}

% ---------------------------------------------------------------------------
% §1 Background — multi-branch agentic execution
% ---------------------------------------------------------------------------
Agentic LLM applications interleave generation with branching.
A coding agent proposes several candidate edits or tool calls for the same repository context, executes one, and discards the rest~\cite{react,sweagent,openhands}; a reasoning harness expands a Tree-of-Thoughts frontier or a self-consistency pool from one partial solution~\cite{tot,selfconsistency,cot}.
Specifically, at each step the serving system starts from a committed context---the \emph{trunk}---and produces $k$ candidate continuations---the \emph{residuals}---that can be independently expanded or evaluated as subsequent requests.
These requests therefore share the trunk but carry distinct residuals; after exploration and selection, at most one path---the \emph{spine}---survives into the next step, while the rest are discarded~\cite{gsm8k,math,esc}.
This multi-branch pattern requires the system to serve a fan-out of candidates even though only one lineage continues.

% ---------------------------------------------------------------------------
% §2 Why important — compute / communication cost under PD disaggregation
% ---------------------------------------------------------------------------
Branching improves search coverage, but it also multiplies serving cost.
A residual that must be expanded further processes any still-uncached input (\emph{prefill}) and then generates tokens (\emph{decode}).
Production stacks increasingly place these phases on separate instances under prefill--decode (PD) disaggregation, so the KV pages produced by prefill must also cross a connector before decode can begin~\cite{distserve,splitwise}.
Most residuals are eventually discarded, yet before they are pruned they may already have consumed prefill compute, KV transfer, and some decode.
As the agent repeats this fan-out over many steps, the cost of non-surviving candidates accumulates.
Stopping an unnecessary residual earlier is therefore a first-class opportunity for efficient PD-disaggregated agentic serving.

% ---------------------------------------------------------------------------
% §3 Existing gap — path pruning vs.\ KV optimization
% ---------------------------------------------------------------------------
Prior work improves efficiency mainly along two axes: path pruning and KV-cache optimization.
Decode-time path pruners cut low-value branches after partial generation: ESC stops sampling once a window agrees~\cite{esc}, Speculative Rejection halts low-reward prefixes at a horizon~\cite{specrej}, DPTS early-stops children after a mini-step~\cite{dpts}, and STOP predicts futility from a short prefix~\cite{stoploss}.
These methods depend on generated text, so by the time they prune, residual prefill and the connector transfer have often already been paid and cannot be recovered.
KV systems attack redundant work on the shared trunk: automatic prefix caching (APC) in vLLM and RadixAttention in SGLang reuse resident trunk pages~\cite{vllm,sglang}, and agent schedulers with KV retention manage context across tool calls~\cite{autellix,continuum,mori,conserve,forkkv}.
Even when the trunk is fully shared, however, each residual's uncached suffix must still be computed and transferred---work that is wasted if that residual is later discarded.
Prefix caching neither rejects a residual nor controls which pages cross the connector.
Combined with PD disaggregation, pruning therefore remains \emph{stage-local}: each stage can discard only work that has already reached it.
The gap is that existing methods reduce duplicate compute and later decode, yet leave the prefill and KV-transfer cost of low-value candidates largely untouched.
We ask how to cut that early-stage cost in PD-disaggregated agent serving without sacrificing reasoning quality.

\begin{figure*}[t]
\centering
\resizebox{0.78\textwidth}{!}{% PD-Prune runtime in DistServe Fig.6 visual language.
% Flow: Requests → Controller → Prefill/Decoding Pruner → Instance.
\begin{tikzpicture}[
  x=1cm, y=1cm,
  font=\sffamily,
  line cap=round, line join=round,
]
\definecolor{dsctrl}{RGB}{233,150,170}
\definecolor{dsmodel}{RGB}{245,215,95}
\definecolor{dsrun}{RGB}{170,220,230}
\definecolor{dsgpu}{RGB}{95,185,115}
\definecolor{dsinst}{RGB}{235,235,235}
\definecolor{dsedge}{RGB}{55,55,55}
\definecolor{dspre}{RGB}{255,200,140}
\definecolor{dsdec}{RGB}{180,210,255}

\path (0,0.75) rectangle (16.2,9.0);

% ---- Requests ----
\foreach \x in {1.35, 2.15, 2.95, 3.75, 4.55} {
  \draw[-{Latex[length=2.4mm,width=1.9mm]}, line width=1.0pt, draw=dsedge]
    (\x,8.75) -- (\x,8.15);
}
\node[font=\footnotesize\sffamily] at (2.95,8.92) {Requests};

% ---- Controller ----
\draw[draw=dsedge, fill=dsctrl, line width=0.8pt]
  (0.70,7.35) rectangle (15.50,8.05);
\node[font=\large\sffamily\bfseries] at (8.10,7.82) {Controller};
\node[font=\scriptsize\sffamily] at (8.10,7.50) {decide before any kernel: residency $m_i^{\Pin},m_i^{\Din}$ and score $s_i$ $\rightarrow$ $a_i\in\{\Hit,\Admit,\Probe,\Reject\}$};

% ---- Controller → Pruners ----
\draw[-{Latex[length=2.4mm,width=1.9mm]}, line width=1.05pt, draw=dsedge]
  (3.15,7.35) -- (3.15,6.85);
\draw[-{Latex[length=2.4mm,width=1.9mm]}, line width=1.05pt, draw=dsedge]
  (13.05,7.35) -- (13.05,6.85);

% ---- Prefill Pruner (left small box) ----
\draw[draw=dsedge, fill=dspre, line width=0.75pt]
  (1.15,6.05) rectangle (5.15,6.75);
\node[font=\small\sffamily\bfseries] at (3.15,6.40) {Prefill Pruner};

% ---- Decoding Pruner (right small box) ----
\draw[draw=dsedge, fill=dsdec, line width=0.75pt]
  (11.05,6.05) rectangle (15.05,6.75);
\node[font=\small\sffamily\bfseries] at (13.05,6.40) {Decoding Pruner};

% ---- Pruners → Instances ----
\draw[-{Latex[length=2.4mm,width=1.9mm]}, line width=1.05pt, draw=dsedge]
  (3.15,6.05) -- (3.15,5.55);
\draw[-{Latex[length=2.4mm,width=1.9mm]}, line width=1.05pt, draw=dsedge]
  (13.05,6.05) -- (13.05,5.55);

% ===================== Prefill Instance (stacked) =====================
\draw[draw=dsedge, fill=dsinst, line width=0.7pt]
  (1.15,0.85) rectangle (6.15,5.05);
\draw[draw=dsedge, fill=dsinst, line width=0.7pt]
  (0.90,1.10) rectangle (5.90,5.30);
\draw[draw=dsedge, fill=dsinst, line width=0.85pt]
  (0.65,1.35) rectangle (5.65,5.55);

\node[font=\normalsize\sffamily\bfseries] at (3.15,5.25) {Prefill Instance};

\draw[draw=dsedge, fill=dsmodel, line width=0.7pt]
  (1.05,4.35) rectangle (5.25,4.95);
\node[font=\small\sffamily] at (3.15,4.65) {LLM Model};

\draw[draw=dsedge, fill=dsrun, line width=0.7pt]
  (1.05,1.65) rectangle (5.25,4.15);
\node[font=\small\sffamily] at (3.15,3.85) {Parallel Runtime};

\foreach \gx/\gy in {1.35/2.85, 3.25/2.85, 1.35/1.90, 3.25/1.90} {
  \draw[draw=dsedge, fill=dsgpu, line width=0.6pt]
    (\gx,\gy) rectangle ({\gx+1.50},{\gy+0.72});
  \node[font=\small\sffamily] at ({\gx+0.75},{\gy+0.36}) {GPU};
}

% ===================== Decoding Instance (stacked) =====================
\draw[draw=dsedge, fill=dsinst, line width=0.7pt]
  (11.05,0.85) rectangle (16.05,5.05);
\draw[draw=dsedge, fill=dsinst, line width=0.7pt]
  (10.80,1.10) rectangle (15.80,5.30);
\draw[draw=dsedge, fill=dsinst, line width=0.85pt]
  (10.55,1.35) rectangle (15.55,5.55);

\node[font=\normalsize\sffamily\bfseries] at (13.05,5.25) {Decoding Instance};

\draw[draw=dsedge, fill=dsmodel, line width=0.7pt]
  (10.95,4.35) rectangle (15.15,4.95);
\node[font=\small\sffamily] at (13.05,4.65) {LLM Model};

\draw[draw=dsedge, fill=dsrun, line width=0.7pt]
  (10.95,1.65) rectangle (15.15,4.15);
\node[font=\small\sffamily] at (13.05,3.85) {Parallel Runtime};

\foreach \gx/\gy in {11.25/2.85, 13.15/2.85, 11.25/1.90, 13.15/1.90} {
  \draw[draw=dsedge, fill=dsgpu, line width=0.6pt]
    (\gx,\gy) rectangle ({\gx+1.50},{\gy+0.72});
  \node[font=\small\sffamily] at ({\gx+0.75},{\gy+0.36}) {GPU};
}

% ===================== Thick hollow KV arrow =====================
\fill[white]
  (5.95,2.85) -- (8.35,2.85) -- (8.35,2.45) -- (9.85,3.20) --
  (8.35,3.95) -- (8.35,3.55) -- (5.95,3.55) -- cycle;
\draw[draw=dsedge, line width=1.25pt]
  (5.95,2.85) -- (8.35,2.85) -- (8.35,2.45) -- (9.85,3.20) --
  (8.35,3.95) -- (8.35,3.55) -- (5.95,3.55) -- cycle;

\node[font=\footnotesize\sffamily\bfseries] at (7.40,4.20) {KV Cache Transfer};
\node[font=\scriptsize\sffamily, align=center] at (7.40,2.05)
  {decoder-missing pages\\[-1pt]of admitted residuals only};

\end{tikzpicture}
}
\caption{\textbf{\sys{} runtime.}
The controller reads the residency of each residual on the prefill instance ($m_i^{\Pin}$) and the decode instance ($m_i^{\Din}$) together with a model-free score $s_i$, and assigns an action $a_i\in\{\Hit,\Admit,\Probe,\Reject\}$ before any kernel runs.
The prefill instance computes only the non-resident suffix of admitted and probed residuals; the connector carries only the pages the decode instance lacks.
The decode-side pruner extends a probed residual beyond $t_0$ tokens only if its generated prefix clears $\tau_{\mathrm{dec}}$.
After commit, only the spine $x_\star$ remains resident.}
\label{fig:system}
\end{figure*}

% ---------------------------------------------------------------------------
% §4 Challenges
% ---------------------------------------------------------------------------
Achieving this goal, however, presents three key challenges.

\paragraph{Challenge 1: Low-overhead pre-execution decision.}
Before a residual runs, the signals that existing pruners rely on---generated tokens---are unavailable, and KV residency is split across instances: a page resident on the prefill side may still require transfer or recompute on the decode side.
The controller must decide, cheaply and from pre-execution information alone, whether each residual should be skipped, prefilled, transferred, or decoded.

\paragraph{Challenge 2: Balancing pruning efficiency and reliability.}
A residual that looks weak before prefill may still reach a correct answer after further generation.
Aggressive pre-kernel rejection saves resources but can hurt quality; retaining every residual preserves quality but erodes the gain.
Candidates of different certainty need different degrees of exploration: a short subsequent generation should refine the decision when the pre-execution signal is inconclusive.

\paragraph{Challenge 3: Pruning-aware KV state management.}
Siblings share the trunk's KV but own distinct residual pages with different lifetimes.
Aborting a branch must free its private pages without dropping pages still referenced by live siblings, while the committed spine must stay resident as the trunk of the next agent step.

% ---------------------------------------------------------------------------
% §5 System + contributions
% ---------------------------------------------------------------------------
To handle these challenges, we design \sys with pre-kernel admission, a decode-side probe, and a copy-on-write context tree (\Cref{fig:system}).
Specifically, (i)~a controller-side admission policy assigns each residual one of $\{\Hit,\Admit,\Probe,\Reject\}$ before any kernel runs, using dual-side page residency and a model-free residual score (\S\ref{sec:app}); (ii)~residuals that cannot be resolved before prefill are \emph{probed}: they are prefilled, decoded for $t_0$ tokens, and extended only if the generated prefix clears a decode-side threshold $\tau_{\mathrm{dec}}$ calibrated separately from the prefill-side threshold $\tau_{\mathrm{pre}}$ (\S\ref{sec:split}), while decode-time hosts such as ESC, Speculative Rejection, and DPTS operate unchanged on the admitted set; and (iii)~a copy-on-write context tree shares the trunk across siblings, transfers only decoder-missing pages, frees residual pages on abort, and retains the committed spine as the trunk of the next agent step (\S\ref{sec:state}).
We call this cross-stage placement \emph{PD pruning disaggregation}: the pruning decision, its execution, and the management of pruned state sit where each is cheapest and most informative.

\paragraph{Results.}
On Qwen3-8B GSM8K with $k{=}16$, admission in front of ESC, Speculative Rejection, or DPTS reduces computed prefill from $5008$ to $3848$ tokens ($-23\%$) for every host; committed-answer accuracy is unchanged, and ESC wall-clock falls from $16.49$\,s to $14.45$\,s (\S\ref{sec:eval}).
On MATH-500, rejecting every low-scoring residual before prefill lowers any-finished coverage from $20/32$ to $13/32$, whereas a $128$-token probe restores $20/32$ at $73.2$\,s instead of $104.1$\,s.
On the full GSM8K test set with Qwen3-14B, the retained spine requires $0.72\times$ the median peak KV of APC at matched accuracy, and the ratio stays within $0.72$--$0.75\times$ across six model families.
All measurements run on physically disaggregated vLLM prefill and decode instances connected by a KV connector.

\paragraph{Contributions.}
Our contributions are summarized as follows:
\begin{itemize}[leftmargin=*,itemsep=1pt,topsep=2pt]
\item We formulate pruning in disaggregated agentic serving as a per-residual, cross-stage decision problem, and make explicit the cache regimes and residency conditions under which prefix caching does or does not share the trunk (\S\ref{sec:problem}).
\item We design \sys, which disaggregates the pruning decision (pre-kernel admission), its execution (prefill, connector, and decode-side probe), and pruned-state management (copy-on-write context tree), and we give a work and transfer accounting that separates trunk sharing, admission, and probe overhead (\S\ref{sec:design}).
\item We show on Qwen3 models that admission composes with three decode-time pruners, that a decode-side probe recovers the coverage lost by blanket pre-kernel rejection, and that spine retention lowers peak KV, while reporting committed-answer accuracy throughout (\S\ref{sec:eval}).
\end{itemize}
